
\documentclass[a4paper,12pt]{article}

\input{glyphtounicode}
\pdfgentounicode=1

\usepackage[T1]{fontenc}
\usepackage[utf8]{inputenc}
\usepackage[english]{babel}
\usepackage[version=4]{mhchem}
\usepackage[dvipsnames]{xcolor}
\usepackage[backend=biber,style=ieee]{biblatex}
\usepackage{newtxtext,newtxmath}
\usepackage{microtype}
\usepackage{setspace}
\usepackage{titlesec}
\usepackage{tocloft}
\usepackage{indentfirst}
\usepackage{enumitem}
\usepackage{fancyhdr}
\usepackage{hyperref}
\usepackage{tabularx}
\usepackage{array}
\usepackage{makecell}
\usepackage{booktabs}
\usepackage{graphicx}
\usepackage{amsmath}
\usepackage{footmisc}
\usepackage{csquotes}
\usepackage{fontawesome5}
\usepackage[
  a4paper,
  top        = 2.00cm,
  bottom     = 2.00cm,
  left       = 2.00cm,
  right      = 2.00cm,
  headheight = 15pt,
  headsep    = 5pt,
  footskip   = 15pt
]{geometry}

\newcommand{\degreeawardvalue}{}
\newcommand{\universityvalue}{}
\newcommand{\addressvalue}{}
\newcommand{\unilogovalue}{}
\newcommand{\copyyearvalue}{}
\newcommand{\defenddatevalue}{}
\newcommand{\orcidvalue}{}
\newcommand{\rightsstatementvalue}{}
\newcommand{\degreeaward}[1]{\def\degreeawardvalue{#1}}
\newcommand{\university}[1]{\def\universityvalue{#1}}
\newcommand{\address}[1]{\def\addressvalue{#1}}
\newcommand{\unilogo}[1]{\def\unilogovalue{#1}}
\newcommand{\copyyear}[1]{\def\copyyearvalue{#1}}
\newcommand{\defenddate}[1]{\def\defenddatevalue{#1}}
\newcommand{\orcid}[1]{\def\orcidvalue{#1}}
\newcommand{\rightsstatement}[1]{\def\rightsstatementvalue{#1}}
\newcommand{\titleskip}{\vskip 4\bigskipamount}
\newcommand{\authorskip}{\vskip 2\bigskipamount}
\newcommand{\resumesection}[2]{\section{\textcolor{red}{#1}#2}}
\newcommand{\resumeSubHeadingListStart}{\begin{itemize}[leftmargin=0.0in, label={}, itemsep=2pt]}
\newcommand{\resumeSubHeadingListEnd}{\end{itemize}}
\newcommand{\resumeItemListStart}{\begin{itemize}}
\newcommand{\resumeItemListEnd}{\end{itemize}}
\newcommand{\resumeItem}[1]{
    \item\small{#1}
}
\newcommand{\resumeSubheading}[4]{
    \item
    \begin{tabular*}{1.0\textwidth}[t]{l@{\extracolsep{\fill}}r}
        \textbf{#1} & \textbf{\small #2} \\
        \small#3 & \small #4
    \end{tabular*}
}

\renewcommand{\cftsecleader}{\cftdotfill{\cftdotsep}}
\renewcommand{\cftsecpresnum}{Chapter~}
\renewcommand{\cftsecaftersnum}{:}
\renewcommand{\headrulewidth}{0pt}
\renewcommand{\footrulewidth}{0pt}
\renewcommand\labelitemi{$\vcenter{\hbox{\tiny$\bullet$}}$}
\renewcommand\labelitemii{$\vcenter{\hbox{\tiny$\bullet$}}$}

\pagestyle{fancy}
\fancyhf{}
\fancyfoot[R]{\thepage}
\urlstyle{same}
\addbibresource{references.bib}
\hypersetup{
  colorlinks=true,
  linkcolor=RoyalBlue,
  urlcolor=RoyalBlue,
  citecolor=RoyalBlue,
  anchorcolor=RoyalBlue
}
\settowidth{\cftsecnumwidth}{Chapter 9: }
\titleformat{\section}
  {\normalfont\Large\bfseries}
  {Chapter \thesection:}
  {0.5em}
  {}
% \titleformat{\section}{
%     \scshape\raggedright\large\bfseries
%     }{}{0em}{}[\color{black}\titlerule]
%     \titlespacing{\section}{0pt}{5pt}{5pt}

\setstretch{1.5}
\setlist[itemize]{topsep=4pt, itemsep=3pt, parsep=0pt, partopsep=0pt, leftmargin=*}
\setlist[enumerate]{topsep=4pt, itemsep=3pt, parsep=0pt, partopsep=0pt, leftmargin=*}
\setlength{\footnotesep}{15pt} 
\setlength{\skip\footins}{15pt} 
\setlength{\parskip}{5pt}

\begin{document}

{\centering
  {\scshape \textls[50]{\fontsize{25}{30}\selectfont Reza {\bfseries Aliasgari Renani}}} \\ 
  {\fontsize{10}{10}\selectfont 31/01/2026} $|$
  {\fontsize{10}{10}\selectfont Motivation Letter} $|$
  {\fontsize{10}{10}\selectfont PhD} $|$
  {\fontsize{10}{10}\selectfont Physics} $|$
  {\fontsize{10}{10}\selectfont PRISMA++} $|$
  {\fontsize{10}{10}\selectfont Johannes Gutenberg University Mainz} $|$
  {\fontsize{10}{10}\selectfont Mu3e experiment} \\
  {\fontsize{10}{10}\selectfont Readout and data transmission electronics for phase II} $|$
  {\fontsize{10}{10}\selectfont Fast online event reconstruction}
\par}
\vspace{-15pt}
\noindent\rule{\textwidth}{0.5pt}

\vspace{-5pt}

My research objective is to develop experimental setups that rely on electronics and reconstruction for particle physics, keeping pace with high rates while preserving timing and resolution. My motivation is to design and program such systems, especially FPGA-based readout and real-time processing, and to connect hardware choices with physics performance. The PRISMA++ PhD Fellowship offers this environment, where detector development, data acquisition, and online analysis are tightly coupled in large international collaborations.

I joined Dr.\ Koveshnikov's laboratory in the Russian Academy of Sciences three years ago to study irradiated semiconductor devices. We investigated the effects of electron and gamma irradiation and hydrogen plasma treatments on \ce{SiO2}-based MOS devices, identifying traps by location (insulator, semiconductor, interface) and carrier type (electrons, holes) using electrical characterization techniques. This work resulted in a first-author publication\footnote{R.~Aliasgari~Renani, O.A.~Soltanovich, M.A.~Knyazev, S.V.~Koveshnikov, ``Investigation of low energy electron irradiated \ce{SiO2} based MOS devices by C--V and TSC techniques,'' \textit{Russian Microelectronics}, 2023, DOI: \href{https://doi.org/10.1134/S1063739723600516}{10.1134/S1063739723600516}} and strengthened my understanding of device stability under bias stress and radiation. I also developed optimized MATLAB automation programs for experimental techniques such as TSC, C--V, I--V, I--T, and DLTS measurements, which enabled systematic data acquisition and reduced experimental run-times.

After my undergraduate degree, I was interested in extending my work from basic semiconductor devices to more complex microelectronic systems and their interaction with radiation. I joined the SoC Development Laboratory to learn about programmable hardware and to work on FPGA-based systems. We implemented image signal processing (ISP) algorithms in Verilog using both manual coding and generated HDL code from Simulink models. I developed a fixed-point arithmetic library, built verification testbenches, and used Python for quality control. We evaluated resource utilization after synthesizing the designs on a Xilinx Zybo Z7 board with live camera input and HDMI output. The codebase was maintained under Git-based source control. 

Concurrently, at the Laboratory of Plasma Systems under the supervision of Dr.\ Vasiliev, we conducted research on radiation and plasma effects on FPGA devices. We exposed an FPGA board to 25--60~keV electron beams in low-pressure oxygen, generating plasma and X-rays from a tungsten target while applying thermal cycling. The FPGA output signal was continuously monitored, and preliminary functionality tests were performed during and after exposure.

I am particularly interested in working under the supervision of Dr.\ Berger on the Mu3e topics of Phase II readout and data transmission electronics (my first choice) and fast online event reconstruction. I plan to contribute to high-rate precision experiments and Mu3e data acquisition and track reconstruction, including tracking methods that remain fast at very high rates and are suitable for GPU processing. This aligns with my experience in FPGA-based DSP, verification, and automated measurement control, which I can apply to firmware and interface testing for the readout chain, systematic performance evaluation, and algorithm--hardware co-development to keep reconstruction efficient as rates scale.

The PhD position at JGU will allow me to bring together my background in characterization of irradiated semiconductor devices and FPGA system development with the demands of a modern high-rate detector. I want to contribute to reliable electronics and online processing that deliver clean detector readout and fast reconstruction of electron tracks and decay vertices in the Mu3e particle-physics experiment, while deepening my expertise in particle detector instrumentation and real-time reconstruction. I am confident that my preparation and motivation fit the PRISMA++ environment and the goals of the Mu3e experiment.

\end{document}
